% EXPERIMENT DRAFT — v2 retrieval study fork (exp-v0.1). NOT part of the v1 paper line.
% Numbers present are from the v1 final_results_bundle.json only. Results section is
% empty by gate discipline until verified probe artifacts exist.
\documentclass{article}
\title{Does Evidence Help? A Retrieval-Augmented Re-Test of Out-of-Context
Detection with a Small Typed-Decision Model}
\author{Rafiul Haider, Pace University (Analytics Capstone, Fall 2026, Team 3)}
\begin{document}
\maketitle
\section{Motivation}
In the v1 study, fine-tuning LAYA-Vision on NewsCLIPpings produced a registered null:
trained accuracy 0.5104 vs a zero-shot baseline of 0.5141 on the full 7,024-row
validation split (H2, not met). The diagnosis localised the missing signal outside
the image--caption pair: out-of-context falsity is a per-item factual property.
This experiment tests whether retrieved evidence supplies that signal.
\section{Design (probe, pre-registered)}
Automated pipeline: caption $\rightarrow$ spaCy entity extraction (deterministic)
with a deepseek-flash claim-formulation variant $\rightarrow$ SearXNG search
$\rightarrow$ timestamped evidence snapshots $\rightarrow$ LAYA judges pair-only vs
pair+evidence on the same 300 validation items (paired comparison, bootstrap
interval). Fact-check domains are excluded in the headline arm and measured in a
second arm. Coverage and abstention are reported, not hidden.

\section{Feature: Search $\rightarrow$ LAYA}
The system feature under test: a SearXNG search is run for each item; the retrieved
evidence is passed into LAYA's context; LAYA then issues the decision. The model is
unchanged --- the variable is the evidence in its context.
\subsection{Attuned questions (registered wording)}
The evidence arm does not reuse the pair-only questions. \textbf{Choice} (native;
pristine/falsified): ``Consider the image, its caption, and the retrieved evidence
below, which reports what is independently known about the people, places, events,
and origins involved. Does the caption accurately describe this image in its
original context? Choose pristine if the caption truthfully describes the image as
originally published, or falsified if the caption misrepresents the image, its
origin, or its context.'' \textbf{Noul} (binary): ``Given the image, the caption,
and the retrieved evidence, what is the probability that this image--caption
pairing is fake --- that the caption does not truthfully describe the image in its
original context?'' The pair-only control arm keeps the v1 pack questions verbatim;
the comparison therefore measures the evidence system as deployed, and the paper
says so.
\subsection{Attuned scoring}
Primary: paired accuracy delta (evidence arm minus pair-only) on identical items,
bootstrap interval. Strata: evidence-found vs no-evidence items reported
separately. Contamination: headline arm excludes fact-check domains; a second arm
includes them; both reported. Arbitration: the evidence-arm choice headlines;
choice/noul disagreement is reported as ``uncertain'' (the v1 rule). Calibration:
evidence-arm temperatures are fitted only on a reserved calibration slice of
train-split items, never on validation, with fit-split and held-out ECE both
reported. Coverage: the share of items with at least one usable evidence passage
inside the token budget is a headline number, not a footnote.

\subsection{Method: baseline first, then a small fine-tune (with train and test)}
Stage A (baseline): the existing v1 NewsCLIPpings checkpoint judges pair+evidence
zero-shot, against the pair-only control on identical items. Stage B (small
fine-tune): a carve from the official train split only (500 train / 100 tune
items), each with SearXNG evidence retrieved once and frozen as snapshots ---
training is reproducible from the snapshots, never from live search, and the
contamination exclusions apply to training evidence as well. The fine-tune uses
the v1 recipe family (full arm, soft\_ce, vision tower frozen, 245 steps) on
examples formatted exactly as the evidence arm is evaluated. Selection on the
tune slice only; evaluation once, on the same validation probe items as Stage A;
the NewsCLIPpings test split stays untouched. Stage A vs Stage B is reported as
a paired comparison on identical items.
\section{Results}
% EMPTY until verified probe artifacts land (see README promotion rule).
\section{Discussion}
% To be written from probe evidence only.
\end{document}
